% Figure: successive integration of a cantilever under uniform load.
% From the distributed load w, one integration gives the shear, a second
% the bending moment, a third the slope, and a fourth the deflection.
% The panels show moment M(x), slope theta(x), and deflection y(x) for a
% cantilever fixed at x = 0 and free at x = L, carrying uniform load w.
\begin{figure}[htbp]
\centering
\begin{tikzpicture}

\begin{axis}[
  name=mom,
  width=5.0cm, height=4.4cm,
  xlabel={$x/L$}, ylabel={$M$},
  xmin=0, xmax=1, ymin=-0.55, ymax=0.05,
  domain=0:1, samples=120,
  axis lines=left, font=\footnotesize,
  title={Bending moment},
]
% M(x) = -w/2 (L-x)^2, normalised by wL^2: -(1/2)(1-x)^2
\addplot[draw=none, fill=engblue!12, domain=0:1]
  {-0.5*(1-x)^2} \closedcycle;
\addplot[engblue, very thick] {-0.5*(1-x)^2};
\end{axis}

\begin{axis}[
  at={(mom.east)}, anchor=west, xshift=14pt,
  name=slp,
  width=5.0cm, height=4.4cm,
  xlabel={$x/L$}, ylabel={$\theta$},
  xmin=0, xmax=1, ymin=-0.02, ymax=0.35,
  domain=0:1, samples=120,
  axis lines=left, font=\footnotesize,
  title={Slope},
]
% theta ~ integral of M: (1/6)(1 - (1-x)^3), normalised
\addplot[engamber, very thick] {(1/6)*(1 - (1-x)^3)};
\end{axis}

\begin{axis}[
  at={(slp.east)}, anchor=west, xshift=14pt,
  width=5.0cm, height=4.4cm,
  xlabel={$x/L$}, ylabel={$y$},
  xmin=0, xmax=1, ymin=-0.02, ymax=0.16,
  domain=0:1, samples=120,
  axis lines=left, font=\footnotesize,
  title={Deflection},
]
% y ~ double integral: (1/24)(6x^2 - 4x^3 + x^4)/... use x^2(6-4x+x^2)/24
\addplot[engred, very thick] {(x^2*(6 - 4*x + x^2))/24};
\end{axis}

\end{tikzpicture}
\caption[Cantilever deflection by successive integration]{A cantilever
fixed at $x = 0$ and free at $x = L$ under a uniform load $w$, worked by
successive integration. The bending moment $M(x) =
-\tfrac{w}{2}(L-x)^{2}$ (left, shaded) is integrated once against the
flexural relation $EI\,y'' = M$ to give the slope $\theta(x)$ (centre)
and a second time to give the deflection $y(x)$ (right), each
integration fixing its constant from a boundary condition at the fixed
end ($y(0) = 0$, $\theta(0) = 0$). The tip deflection $y(L) =
wL^{4}/(8EI)$ falls out as the value of the second integral at the free
end. Each curve is the running integral of the one to its left, the same
accumulation the chapter's fundamental theorem describes, applied twice
in series.}
\label{fig:vol02ch06:beam-deflection}
\end{figure}
